% Generated by tools/docs/gen_fact_tables.py from rtl/soc/pmu.sv and docs/manual/common/facts/power.yaml. Do not edit.
{\small
\begin{longtable}{@{}>{\raggedright\arraybackslash}p{31mm}>{\raggedright\arraybackslash}p{9mm}>{\raggedright\arraybackslash}p{19mm}>{\raggedright\arraybackslash}p{22mm}>{\raggedright\arraybackslash}p{58mm}@{}}
\caption{Signals of \code{pmu}}\label{tab:tm:mod:pmu:signals}\\
\toprule
Signal & Dir & Width & Clock domain & Description\\
\midrule
\endfirsthead
\multicolumn{5}{@{}l}{\emph{Table \thetable\ continued}}\\
\toprule
Signal & Dir & Width & Clock domain & Description\\
\midrule
\endhead
\bottomrule
\endlastfoot
\code{clk_i} & in & 1 & clock & Clock of the unit. In the SoC it is \code{core_clk}.\\
\code{rst_n_i} & in & 1 & \code{clk_i} & Active-low synchronous reset. In the SoC it is \code{core_rst_n}.\\
\code{psel} & in & 1 & \code{clk_i} & APB select.\\
\code{penable} & in & 1 & \code{clk_i} & APB enable. It is high in the ACCESS phase.\\
\code{pwrite} & in & 1 & \code{clk_i} & APB direction. High for a write.\\
\code{paddr} & in & \code{ADDR_W} (12) & \code{clk_i} & APB byte address inside the slot. Bits 1 and 0 are not used.\\
\code{pwdata} & in & 32 & \code{clk_i} & APB write data.\\
\code{pstrb} & in & 4 & \code{clk_i} & APB byte strobes.\\
\code{prdata} & out & 32 & \code{clk_i} & APB read data. It is a combinational function of the address and the registers.\\
\code{pready} & out & 1 & \code{clk_i} & APB ready. It is tied high. The unit adds no wait state.\\
\code{pslverr} & out & 1 & \code{clk_i} & APB error. It is tied low. The unit never gives an error.\\
\code{cpu_clk_en_o} & out & 1 & \code{clk_i} & Clock enable of the CPU domain. High means the clock runs. The output is the Q of a flop.\\
\code{cpu_iso_en_o} & out & 1 & \code{clk_i} & Isolation control of the CPU domain. High means the clamp is asked. The output is the Q of a flop.\\
\code{cpu_rst_n_o} & out & 1 & \code{clk_i} & Active-low reset of the CPU domain. The output is the Q of a flop.\\
\code{cpu_ret_save_o} & out & 1 & \code{clk_i} & Retention save pulse of the CPU domain. It is high for one cycle in each power-down. No retention cell uses it.\\
\code{cpu_ret_restore_o} & out & 1 & \code{clk_i} & Retention restore pulse of the CPU domain. It is high for one cycle in each power-up. No retention cell uses it.\\
\code{cpu_iso_ack_i} & in & 1 & \code{clk_i} & Isolation acknowledge of the CPU domain, already synchronised into \code{clk_i} by the caller. High means the clamp is seen active. Low means the clamp is seen released.\\
\code{gpu_clk_en_o} & out & 1 & \code{clk_i} & Clock enable of the GPU domain. High means the clock runs. The output is the Q of a flop.\\
\code{gpu_iso_en_o} & out & 1 & \code{clk_i} & Isolation control of the GPU domain. High means the clamp is active. The output is the Q of a flop.\\
\code{gpu_rst_n_o} & out & 1 & \code{clk_i} & Active-low reset of the GPU domain. The output is the Q of a flop.\\
\code{gpu_ret_save_o} & out & 1 & \code{clk_i} & Retention save pulse of the GPU domain. It is high for one cycle in each power-down. No retention cell uses it.\\
\code{gpu_ret_restore_o} & out & 1 & \code{clk_i} & Retention restore pulse of the GPU domain. It is high for one cycle in each power-up. No retention cell uses it.\\
\end{longtable}
}
